Las notaciones asintóticas permiten clasificar la eficiencia de los algoritmos acotando su tiempo de ejecución y requerimiento de memoria en función del tamaño de entrada. Sin embargo, estos modelos teóricos simplifican el comportamiento práctico al omitir las constantes multiplicativas. En este trabajo se analiza experimentalmente el desempeño de cuatro algoritmos de ordenamiento (\textsc{sort}, \textsc{QuickSort}, \textsc{MergeSort} y \textsc{PatienceSort}) y dos de multiplicación de matrices cuadradas (\textsc{Naive} y \textsc{Strassen}). El objetivo consiste en contrastar sus cotas teóricas conocidas con los resultados obtenidos de tiempo y memoria frente a variaciones en la disposición de los datos, el dominio numérico y la dimensión de las entradas.